% !TeX program = xelatex
% UTF-8; compile twice with XeLaTeX. No external images or data files needed.
\documentclass[a4paper,landscape,10pt]{article}
\usepackage[margin=14mm,headheight=14pt,headsep=5mm,footskip=7mm]{geometry}
\usepackage{lmodern}
\usepackage{fontspec}
\usepackage{url}
\setmainfont{DejaVu Sans}
\newfontfamily\ipafont{DejaVu Sans}
\usepackage[table]{xcolor}
\usepackage{array,longtable,booktabs,fancyhdr}
\definecolor{ink}{HTML}{18384D}
\definecolor{band}{HTML}{E7EFF4}
\pagestyle{fancy}\fancyhf{}
\fancyhead[L]{\small\color{ink}Parlez Vous QSO?! · Česky / Czech}
\fancyhead[R]{\small Englische Phrase · Český překlad · IPA}
\fancyfoot[L]{\scriptsize Vorlage: Czech v 1.1 · Credits: Josef -- OK1UNE}
\fancyfoot[R]{\small\thepage}
\setlength{\parindent}{0pt}
\setlength{\tabcolsep}{3mm}
\renewcommand{\arraystretch}{1.25}
\newcommand{\ipa}[1]{{\ipafont[#1]}}
\begin{document}
{\LARGE\bfseries\color{ink}Parlez Vous QSO?!}\hfill{\large Česky / Czech}

\vspace{2mm}
{\small QSO-Sprachführer mit tschechischer IPA-Lautschrift. Die Lautschrift bezieht sich auf die tschechische Spalte.}

\vspace{2mm}
{\footnotesize \textbf{Leseschlüssel:} \ipa{ˈ} Betonung; \ipa{ː} langer Vokal (unabhängig von der Betonung);
\ipa{r̝} stimmhaftes ř, \ipa{r̝̊} stimmloses ř; \ipa{c, ɟ, ɲ} entsprechen ť, ď, ň;
\ipa{x} wie „ch“ in „Bach“. \ipa{r̩, l̩} sind silbische Konsonanten.}

\vspace{1mm}
\fontsize{9}{11.8}\selectfont
\begin{longtable}{@{}>{\raggedright\arraybackslash}p{13mm}>{\raggedright\arraybackslash}p{73mm}>{\raggedright\arraybackslash}p{75mm}>{\raggedright\arraybackslash}p{90mm}@{}}
\toprule
\textbf{Ref.} & \textbf{English phrase} & \textbf{Česky} & \textbf{Lautschrift (IPA)}\\
\midrule
\endfirsthead
\toprule
\textbf{Ref.} & \textbf{English phrase} & \textbf{Česky} & \textbf{Lautschrift (IPA)}\\
\midrule
\endhead
\midrule
\multicolumn{4}{r}{\scriptsize Fortsetzung auf der nächsten Seite}\\
\endfoot
\bottomrule
\endlastfoot
\rowcolor{band} 1.0 & \textbf{Opening Greetings} & \multicolumn{2}{l}{\textbf{Úvodní pozdravy}} \\*
1.1 & I am very pleased to meet you. & Velmi mě těší, že vás poznávám. & \ipa{ˈvɛlmɪ mɲɛ ˈcɛʃiː ʒɛ vaːs ˈpoznaːvaːm} \\
1.2 & Thank you for your call & Děkuji vám za zavolání. & \ipa{ˈɟɛkujɪ vaːm za ˈzavolaːɲiː} \\
1.3 & Good morning & Dobré ráno. & \ipa{ˈdobrɛː ˈraːno} \\
1.31 & Good day & Dobrý den. & \ipa{ˈdobriː ˈdɛn} \\
1.32 & Good afternoon & Dobré odpoledne. & \ipa{ˈdobrɛː ˈotpolɛdnɛ} \\
1.33 & Good evening & Dobrý večer. & \ipa{ˈdobriː ˈvɛtʃɛr} \\
1.4 & I know only a little Czech but I enjoy QSOs in Czech. & Znám jenom trochu česky, ale mám rád spojení v češtině. & \ipa{znaːm ˈjɛnom ˈtroxu ˈtʃɛskɪ ˈalɛ maːm raːt ˈspojɛɲiː f ˈtʃɛʃcɪɲɛ} \\
\rowcolor{band} 2.0 & \textbf{Operator} & \multicolumn{2}{l}{\textbf{Operátor}} \\*
2.1 & My name is Peter. & Jmenuji se Peter. & \ipa{ˈjmɛnujɪ sɛ ˈpiːtə} \\
2.2 & I have been licensed since 1960. & Mám licenci od roku 1960. & \ipa{maːm ˈlɪtsɛntsɪ ot ˈroku ˈcɪsiːts ˈɟɛvɛt sɛt ˈʃɛdɛsaːt} \\
2.3 & I have been licensed for 43 years. & Mám licenci 43 let. & \ipa{maːm ˈlɪtsɛntsɪ ˈtʃtɪr̝ɪtsɛt ˈtr̝̊ɪ lɛt} \\
2.4 & I live in Sittingbourne. & Bydlím v Sittingbourne. & \ipa{ˈbɪdliːm f ˈsɪtɪŋbɔːn} \\
2.5 & Sittingbourne is a village. & Sittingbourne je vesnice. & \ipa{ˈsɪtɪŋbɔːn jɛ ˈvɛsɲɪtsɛ} \\
2.51 & Sittingbourne is a town. & Sittingbourne je městečko. & \ipa{ˈsɪtɪŋbɔːn jɛ ˈmɲɛstɛtʃko} \\
2.52 & Sittingbourne is a city. & Sittingbourne je velké město. & \ipa{ˈsɪtɪŋbɔːn jɛ ˈvɛlkɛː ˈmɲɛsto} \\
2.6 & Sittingbourne is between London and Dover. & Sittingbourne je situováno mezi Londýnem a Doverem. & \ipa{ˈsɪtɪŋbɔːn jɛ ˈsɪtuovaːno ˈmɛzɪ ˈlondiːnɛm a ˈdovɛrɛm} \\
2.7 & I live in the South of England. & Bydlím na jihu Anglie. & \ipa{ˈbɪdliːm ˈna jɪɦu ˈaŋɡlɪjɛ} \\
2.71 & I live in the North of England. & Bydlím na severu Anglie. & \ipa{ˈbɪdliːm ˈna sɛvɛru ˈaŋɡlɪjɛ} \\
2.72 & I live in the East of England. & Bydlím na východu Anglie. & \ipa{ˈbɪdliːm ˈna viːxodu ˈaŋɡlɪjɛ} \\
2.73 & I live in West of England. & Bydlím na západu Anglie. & \ipa{ˈbɪdliːm ˈna zaːpadu ˈaŋɡlɪjɛ} \\
2.74 & I live in South East England. & Bydlím na jiho-východu Anglie. & \ipa{ˈbɪdliːm ˈna jɪɦoviːxodu ˈaŋɡlɪjɛ} \\
2.75 & I live in South West England. & Bydlím na jiho-západu Anglie. & \ipa{ˈbɪdliːm ˈna jɪɦozaːpadu ˈaŋɡlɪjɛ} \\
2.76 & I live in North East England. & Bydlím na severo-východu Anglie. & \ipa{ˈbɪdliːm ˈna sɛvɛroviːxodu ˈaŋɡlɪjɛ} \\
2.77 & I live in North West England. & Bydlím na severo-západu Anglie. & \ipa{ˈbɪdliːm ˈna sɛvɛrozaːpadu ˈaŋɡlɪjɛ} \\
2.78 & I live in central England. & Bydlím ve střední Anglii. & \ipa{ˈbɪdliːm ˈvɛ str̝̊ɛdɲiː ˈaŋɡlɪjɪ} \\
2.8 & I live in the country. & Bydlím na venkově. & \ipa{ˈbɪdliːm ˈna vɛŋkovjɛ} \\
2.81 & I live by the sea. & Bydlím u moře. & \ipa{ˈbɪdliːm ˈu mor̝ɛ} \\
2.82 & I live near mountains. & Bydlím v horách. & \ipa{ˈbɪdliːm ˈv ɦoraːx} \\
2.9 & Sittingbourne is 50Km from London. & Sittingbourne je 50 km od Londýna. & \ipa{ˈsɪtɪŋbɔːn jɛ ˈpadɛsaːt ˈkɪlomɛtruː od ˈlondiːna} \\
2.11 & I am retired. & Jsem důchodce. & \ipa{sɛm ˈduːxottsɛ} \\
2.12 & My home call sign is M0ECS. & Má domovská volací značka je M0ECS. & \ipa{maː ˈdomofskaː ˈvolatsiː ˈznatʃka jɛ ɛm ˈnula ɛː tsɛː ɛs} \\
\rowcolor{band} 3.0 & \textbf{Station} & \multicolumn{2}{l}{\textbf{Stanice}} \\*
3.1 & My radio is an Icom 706. & Moje zařízení je Icom 706. & \ipa{ˈmojɛ ˈzar̝iːzɛɲiː jɛ ˈaɪkɒm ˈsɛdm̩ sɛt ˈʃɛst} \\
3.2 & My station is all homebrew (homemade). & Moji stanici jsem si zhotovil sám. & \ipa{ˈmojɪ ˈstaɲɪtsɪ sɛm sɪ ˈzɦotovɪl saːm} \\
3.3 & I am running 5 Watts. & Pracuji na 5 wattech. & \ipa{ˈpratsujɪ ˈna pjɛtɪ ˈvatɛx} \\
3.4 & My antenna is a long wire. & Mám drátovou anténu. & \ipa{maːm ˈdraːtovou̯ ˈantɛːnu} \\
3.41 & My antenna is a dipole. & Mám dipól. & \ipa{maːm ˈdɪpoːl} \\
3.42 & My antenna is a G5RV. & Mám anténu G5RV. & \ipa{maːm ˈantɛːnu ɡɛː ˈpjɛt ɛr vɛː} \\
3.43 & My antenna is a doublet. & Mám dublet. & \ipa{maːm ˈdublɛt} \\
3.44 & My antenna is a delta loop. & Mám delta loop. & \ipa{maːm ˈdɛlta luːp} \\
3.45 & My antenna is a vertical. & Mám vertikál. & \ipa{maːm ˈvɛrcɪkaːl} \\
3.46 & My antenna is a log periodic. & Mám logaritmicko-periodickou anténu. & \ipa{maːm ˈloɡarɪtmɪtsko pɛrɪodɪtskou̯ ˈantɛːnu} \\
3.47 & My antenna is a Yagi beam. & Mám směrovou anténu Yagi. & \ipa{maːm ˈsmɲɛrovou̯ ˈantɛːnu ˈjaɡɪ} \\
3.48 & My antenna is a 5 element beam. & Mám pětiprvkovou směrovou anténu. & \ipa{maːm ˈpjɛcɪpr̩fkovou̯ ˈsmɲɛrovou̯ ˈantɛːnu} \\
\rowcolor{band} 4.0 & \textbf{Conditions} & \multicolumn{2}{l}{\textbf{Podmínky}} \\*
4.1 & Thank you for my 599 report. & Děkuji vám za report 599. & \ipa{ˈɟɛkujɪ vaːm za ˈrɛport ˈpjɛt ˈɟɛvɛt ˈɟɛvɛt} \\
4.2 & The band conditions are good. & Podmínky na pásmu jsou dobré. & \ipa{ˈpodmiːŋkɪ ˈna paːsmu sou̯ ˈdobrɛː} \\
4.3 & The band conditions are bad. & Podmínky na pásmu jsou špatné. & \ipa{ˈpodmiːŋkɪ ˈna paːsmu sou̯ ˈʃpatnɛː} \\
4.4 & There is strong interference here. & Je zde silné rušení - interference. & \ipa{jɛ zdɛ ˈsɪlnɛː ˈruʃɛɲiː / ˈɪntɛrfɛrɛntsɛ} \\
4.5 & Can you change to the 80m band? & Můžete přejít na pásmo 80m? & \ipa{ˈmuːʒɛtɛ ˈpr̝̊ɛjiːt ˈna paːsmo ˈosm̩dɛsaːt ˈmɛtruː} \\
4.6 & Please change frequency to 3.558 MHz. & Změňte prosím frekvenci na 3.558 MHz. & \ipa{ˈzmɲɛɲtɛ ˈprosiːm ˈfrɛkvɛntsɪ na ˈtr̝̊ɪ ˈtsɛlaː ˈpjɛt ˈpjɛt ˈosm̩ ˈmɛɡaɦɛrtsuː} \\
\rowcolor{band} 5.0 & \textbf{General Queries and Comments} & \multicolumn{2}{l}{\textbf{Základní Otázky a Připomínky}} \\*
5.1 & What does <\#\#\#\#\#> mean please? & Co <\#\#\#\#\#> prosím znamená? & \ipa{tso … ˈprosiːm ˈznamɛnaː} \\
5.2 & Please wait a moment, I need to check my Czech dictionary. & Prosím, počkejte chvíli, potřebuji se podívat do českého slovníku. & \ipa{ˈprosiːm ˈpotʃkɛjtɛ ˈxviːlɪ ˈpotr̝̊ɛbujɪ sɛ ˈpoɟiːvat ˈdo tʃɛskɛːɦo ˈslovɲiːku} \\
5.3 & How do you say, <\#\#\#\#\#> in Czech ? & Jak se řekne <\#\#\#\#\#> v češtině? & \ipa{jak sɛ ˈr̝ɛknɛ … f ˈtʃɛʃcɪɲɛ} \\
5.5 & My favourite modes of operation are CW / RTTY / PSK / MT63 / Olivia / Hellschreiber / Digital Modes / and SSB & Moje oblíbené režimy provozu jsou CW / RTTY / PSK / MT63 / Olivia / Hellschreiber / digitální režimy / a SSB. & \ipa{ˈmojɛ ˈobliːbɛnɛː ˈrɛʒɪmɪ ˈprovozu sou̯ tsɛː ˈdvojɪtɛː vɛː / ɛr tɛː tɛː ˈɪpsɪlon / pɛː ɛs kaː / ɛm tɛː ˈʃɛdɛsaːt ˈtr̝̊ɪ / ˈolɪvɪja / ˈhɛlˌʃʁaɪbɐ / ˈdɪɡɪtaːlɲiː ˈrɛʒɪmɪ / a ɛs ɛs bɛː} \\
5.6 & My best DX is Australia. & Můj nejlepší DX je Austrálie. & \ipa{muːj ˈnɛjlɛpʃiː dɛː ɪks jɛ ˈau̯straːlɪjɛ} \\
5.7 & I enjoy working contests. & Líbí se mi závody. & \ipa{ˈliːbiː sɛ mɪ ˈzaːvodɪ} \\
5.8 & I like working QRP. & Rád používám QRP. & \ipa{raːt ˈpo.uʒiːvaːm kvɛː ɛr pɛː} \\
\rowcolor{band} 6.0 & \textbf{The Weather} & \multicolumn{2}{l}{\textbf{Počasí}} \\*
6.1 & It is sunny. & Je slunečno. & \ipa{jɛ ˈslunɛtʃno} \\
6.2 & It is windy. & Je silný vítr. & \ipa{jɛ ˈsɪlniː ˈviːtr̩} \\
6.3 & It is raining. & Prší. & \ipa{ˈpr̩ʃiː} \\
6.4 & It is snowing. & Sněží. & \ipa{ˈsɲɛʒiː} \\
6.5 & It is very hot. & Je velmi horko. & \ipa{jɛ ˈvɛlmɪ ˈɦorko} \\
6.6 & It is very cold. & Je velká zima. & \ipa{jɛ ˈvɛlkaː ˈzɪma} \\
6.7 & The temperature is 8C. & Teplota je 8 stupňů Celsia. & \ipa{ˈtɛplota jɛ ˈosm̩ ˈstupɲuː ˈtsɛlzɪja} \\
6.8 & It is overcast. & Je zataženo. & \ipa{jɛ ˈzataʒɛno} \\
6.9 & It is foggy. & Je mlha. & \ipa{jɛ ˈml̩ɦa} \\
6.11 & It has been raining all day. & Dnes prší celý den. & \ipa{dnɛs ˈpr̩ʃiː ˈtsɛliː dɛn} \\
6.12 & There is 20cm of snow here. & Je zde 20 cm sněhu. & \ipa{jɛ zdɛ ˈdvatsɛt ˈtsɛncɪmɛtruː ˈsɲɛɦu} \\
\rowcolor{band} 8.0 & \textbf{QSL} & \multicolumn{2}{l}{\textbf{QSL}} \\*
8.1 & My QSL card is sure via the bureau & Můj QSL lístek půjde přes QSL službu. & \ipa{muːj kvɛː ɛs ɛl ˈliːstɛk ˈpuːjdɛ pr̝̊ɛs kvɛː ɛs ɛl ˈsluʒbu} \\
8.2 & QSL via eQSL & QSL jako eQSL. & \ipa{kvɛː ɛs ɛl ˈjako iː kvɛː ɛs ɛl} \\
8.3 & No eQSL please & Prosím ne eQSL. & \ipa{ˈprosiːm nɛ iː kvɛː ɛs ɛl} \\
8.4 & I will send you my card direct. & Pošlu vám můj QSL lístek přímo. & \ipa{ˈpoʃlu vaːm muːj kvɛː ɛs ɛl ˈliːstɛk ˈpr̝̊iːmo} \\
8.5 & What is your address? & Jaká je vaše adresa? & \ipa{ˈjakaː jɛ ˈvaʃɛ ˈadrɛsa} \\
8.6 & Please send your QSL card direct. & Prosím pošlete mi svůj QSL lístek přímo. & \ipa{ˈprosiːm ˈpoʃlɛtɛ mɪ svuːj kvɛː ɛs ɛl ˈliːstɛk ˈpr̝̊iːmo} \\
8.7 & My address is... & Moje adresa je … & \ipa{ˈmojɛ ˈadrɛsa jɛ …} \\
8.8 & My address is shown on QSL.com / EHam.net & Moji adresu naleznete na QSL.com / EHam.net. & \ipa{ˈmojɪ ˈadrɛsu ˈnalɛznɛtɛ na kvɛː ɛs ɛl ˈtɛtʃka kom / iː hæm ˈtɛtʃka nɛt} \\
8.9 & Sorry I do not send QSL cards. & Lituji, QSL lístky neposílám. & \ipa{ˈlɪtujɪ kvɛː ɛs ɛl ˈliːstkɪ ˈnɛposiːlaːm} \\
8.11 & Please send QSL card via G1JQH. & Prosím, pošlete QSL lístek přes G1JQH. & \ipa{ˈprosiːm ˈpoʃlɛtɛ kvɛː ɛs ɛl ˈliːstɛk pr̝̊ɛz ɡɛː ˈjɛdna jɛː kvɛː ɦaː} \\
8.12 & My WACRAL / GQRP / DXCC / FISTS number is 1234. & Moje WACRAL / GQRP / DXCC / FISTS číslo je 1234. & \ipa{ˈmojɛ ˈdvojɪtɛː vɛː aː tsɛː ɛr ɛl / ɡɛː kvɛː ɛr pɛː / dɛː ɪks tsɛː tsɛː / fɪsts ˈtʃiːslo jɛ ˈcɪsiːts ˈdvjɛ scɛ ˈtr̝̊ɪtsɛt ˈtʃtɪr̝ɪ} \\
\rowcolor{band} 9.0 & \textbf{Final Greetings} & \multicolumn{2}{l}{\textbf{Závěrečné pozdravy}} \\*
9.1 & Thank you. & Děkuji vám. & \ipa{ˈɟɛkujɪ vaːm} \\
9.2 & Thank you for a very enjoyable contact. & Děkuji vám za velmi příjemné spojení. & \ipa{ˈɟɛkujɪ vaːm za ˈvɛlmɪ ˈpr̝̊iːjɛmnɛː ˈspojɛɲiː} \\
9.3 & I enjoyed our chat. & Naše povídání se mi líbilo. & \ipa{ˈnaʃɛ ˈpoviːdaːɲiː sɛ mɪ ˈliːbɪlo} \\
9.4 & I enjoyed our QSO in Czech. & Líbilo se mi naše spojení v češtině. & \ipa{ˈliːbɪlo sɛ mɪ ˈnaʃɛ ˈspojɛɲiː f ˈtʃɛʃcɪɲɛ} \\
9.5 & It has been great fun. & Bylo to velmi pěkné. & \ipa{ˈbɪlo to ˈvɛlmɪ ˈpjɛknɛː} \\
9.6 & If you would like a regular contact with me, you can contact me via email. & Pokud chcete být se mnou v pravidelném spojení, můžete mě kontaktovat přes e-mail. & \ipa{ˈpokut ˈxtsɛtɛ biːt sɛ mnou̯ v ˈpravɪdɛlnɛːm ˈspojɛɲiː ˈmuːʒɛtɛ mɲɛ ˈkontaktovat pr̝̊ɛs ˈiːmɛjl} \\
9.7 & My email address is... & Moje e-mailová adresa je … & \ipa{ˈmojɛ ˈiːmɛjlovaː ˈadrɛsa jɛ …} \\
9.8 & My email address is shown on QSL.com / EHam.net & Moji e-mailovou adresu naleznete na QSL.com / EHam.net & \ipa{ˈmojɪ ˈiːmɛjlovou̯ ˈadrɛsu ˈnalɛznɛtɛ na kvɛː ɛs ɛl ˈtɛtʃka kom / iː hæm ˈtɛtʃka nɛt} \\
9.9 & I look forward to meeting you again. & Těším se, až se opět setkáme. & \ipa{ˈcɛʃiːm sɛ aʃ sɛ ˈopjɛt ˈsɛtkaːmɛ} \\
9.11 & My best wishes to you and your family. & Srdečný pozdrav vám a vaší rodině. & \ipa{ˈsr̩dɛtʃniː ˈpozdraf vaːm a ˈvaʃiː ˈroɟɪɲɛ} \\
9.12 & Best wishes to your wife. & Srdečný pozdrav vaší paní. & \ipa{ˈsr̩dɛtʃniː ˈpozdraf ˈvaʃiː ˈpaɲiː} \\
9.121 & Best wishes to your husband. & Srdečný pozdrav vašemu manželovi. & \ipa{ˈsr̩dɛtʃniː ˈpozdraf ˈvaʃɛmu ˈmanʒɛlovɪ} \\
9.13 & Good luck. & Hodně štěstí. & \ipa{ˈɦoɟɲɛ ˈʃcɛsciː} \\
9.14 & Good health. & Pevné zdraví. & \ipa{ˈpɛvnɛː ˈzdraviː} \\
9.15 & Goodbye & Nashledanou. & \ipa{ˈnasxlɛdanou̯} \\
9.16 & I wish you an enjoyable afternoon. & Přeji vám pěkné odpoledne. & \ipa{ˈpr̝̊ɛjɪ vaːm ˈpjɛknɛː ˈotpolɛdnɛ} \\
9.161 & I wish you an enjoyable evening. & Přeji vám pěkný večer. & \ipa{ˈpr̝̊ɛjɪ vaːm ˈpjɛkniː ˈvɛtʃɛr} \\
9.162 & I wish you an enjoyable weekend. & Přeji vám pěkný víkend. & \ipa{ˈpr̝̊ɛjɪ vaːm ˈpjɛkniː ˈviːkɛnt} \\
9.17 & Enjoy your holiday. & Krásnou dovolenou. & \ipa{ˈkraːsnou̯ ˈdovolɛnou̯} \\
9.18 & Happy New Year. & Šťastný nový rok. & \ipa{ˈʃcasniː ˈnoviː rok} \\
\rowcolor{band} 10.0 & \textbf{Country} & \multicolumn{2}{l}{\textbf{Země}} \\*
10.1 & Albania & Albánie & \ipa{ˈalbaːɲɪjɛ} \\
10.2 & Austria & Rakousko & \ipa{ˈrakou̯sko} \\
10.3 & Belgium & Belgie & \ipa{ˈbɛlɡɪjɛ} \\
10.4 & Belorussia & Bělorusko & \ipa{ˈbjɛlorusko} \\
10.5 & Bosnia \& Herzegovina & Bosna a Hercegovina & \ipa{ˈbosna a ˈɦɛrtsɛɡovɪna} \\
10.6 & Bulgaria & Bulharsko & \ipa{ˈbulɦarsko} \\
10.7 & Croatia & Chorvatsko & \ipa{ˈxorvatsko} \\
10.8 & Czech Republic & Česká republika & \ipa{ˈtʃɛskaː ˈrɛpublɪka} \\
10.9 & Denmark & Dánsko & \ipa{ˈdaːnsko} \\
10.11 & Eire & Irská republika & \ipa{ˈɪrskaː ˈrɛpublɪka} \\
10.12 & England & Anglie & \ipa{ˈaŋɡlɪjɛ} \\
10.13 & Estonia & Estonsko & \ipa{ˈɛstonsko} \\
10.14 & Finland & Finsko & \ipa{ˈfɪnsko} \\
10.15 & France & Francie & \ipa{ˈfrantsɪjɛ} \\
10.16 & Germany & Německo & \ipa{ˈɲɛmɛtsko} \\
10.17 & Gibraltar & Gibraltar & \ipa{ˈɡɪbraltar} \\
10.18 & Greece & Řecko & \ipa{ˈr̝ɛtsko} \\
10.19 & Holland & Holandsko & \ipa{ˈɦolantsko} \\
10.21 & Hungary & Maďarsko & \ipa{ˈmaɟarsko} \\
10.22 & Iceland & Island & \ipa{ˈɪslant} \\
10.23 & Northern Ireland & Severní Irsko & \ipa{ˈsɛvɛrɲiː ˈɪrsko} \\
10.24 & Isle of Man & Ostrov Man & \ipa{ˈostrof ˈman} \\
10.25 & Italy & Itálie & \ipa{ˈɪtaːlɪjɛ} \\
10.26 & Latvia & Lotyšsko & \ipa{ˈlotɪʃsko} \\
10.27 & Liechtenstein & Liechtenstein & \ipa{ˈliçtn̩ʃtaɪn} \\
10.28 & Lithuania & Litva & \ipa{ˈlɪtva} \\
10.29 & Luxembourg & Luxemburg & \ipa{ˈluksɛmburk} \\
10.31 & Macedonia & Makedonie & \ipa{ˈmakɛdoɲɪjɛ} \\
10.32 & Moldavia & Moldavsko & \ipa{ˈmoldafsko} \\
10.33 & Netherlands & Holandsko & \ipa{ˈɦolantsko} \\
10.34 & Norway & Norsko & \ipa{ˈnorsko} \\
10.35 & Poland & Polsko & \ipa{ˈpolsko} \\
10.36 & Portugal & Portugalsko & \ipa{ˈportuɡalsko} \\
10.37 & Romania & Rumunsko & \ipa{ˈrumunsko} \\
10.38 & Russian Federation & Ruská federace & \ipa{ˈruskaː ˈfɛdɛratsɛ} \\
10.39 & Scotland & Skotsko & \ipa{ˈskotsko} \\
10.41 & Slovakia & Slovensko & \ipa{ˈslovɛnsko} \\
10.42 & Slovenia & Slovinsko & \ipa{ˈslovɪnsko} \\
10.43 & Spain & Španělsko & \ipa{ˈʃpaɲɛlsko} \\
10.44 & Svalbard & Špicberky (Svalbard) & \ipa{ˈʃpɪtsbɛrkɪ (ˈsvalbart)} \\
10.45 & Sweden & Švédsko & \ipa{ˈʃvɛːtsko} \\
10.46 & Switzerland & Švýcarsko & \ipa{ˈʃviːtsarsko} \\
10.47 & Turkey & Turecko & \ipa{ˈturɛtsko} \\
10.48 & Ukraine & Ukrajina & \ipa{ˈukrajɪna} \\
10.49 & Wales & Wales & \ipa{ˈvɛjls} \\
\end{longtable}
\normalsize
\section*{Hinweise zur Übertragung und Aussprache}
\fontsize{8.5}{10.5}\selectfont\raggedright
\textbf{Vorlage und Umfang.} Grundlage ist „Czech v 1.1.doc“; Credit: Josef -- OK1UNE.
Alle 142 Phrasen und Ländernamen sind enthalten. Die ursprüngliche Gliederung mit ihren Nummernsprüngen
und dem fehlenden Abschnitt 7 bleibt erhalten.

\smallskip
\textbf{Dokumentierte Korrekturen und Ergänzungen.} Die Referenz „.3“ wurde zu „8.3“ berichtigt.
„Watech“ wurde zu „wattech“ (3.3), „5-ti prvkovou“ zu „pětiprvkovou“ (3.48) und „vitr“ zu „vítr“ (6.2).
Die grammatischen Formen „pěkné večer“ und „pěkné víkend“ wurden zu „pěkný večer“ (9.161) und
„pěkný víkend“ (9.162). „Nový rok“ wurde in 9.18 als „nový rok“ geschrieben, weil der Wunsch
für das neue Jahr gilt. Mehrfache Leerzeichen wurden vereinheitlicht. Bei Svalbard (10.44) stand in der
Quelle nur „~“; ergänzt wurde „Špicberky (Svalbard)“. Diese Bezeichnung verwendet auch das tschechische
Außenministerium (Quelle unten). Die Lautschrift folgt jeweils dem korrigierten beziehungsweise ergänzten Text.

\smallskip
\textbf{Inhaltliche Eigenheiten.} Die Übersetzungen bleiben ansonsten wie in der Quelle, etwa
„near mountains“ / „v horách“ (2.82) und „It has been raining all day“ / „Dnes prší celý den“ (6.11).
„Holandsko“ steht bei Holland und Netherlands (10.19 und 10.33) jeweils wie in der Vorlage.
Historische Ländernamen und Webadressen wurden nicht modernisiert. Die ursprünglichen Bindestriche
in den Himmelsrichtungen sowie „logaritmicko-periodickou“ bleiben erhalten. „Rád“ in 1.4 und 5.8
steht für einen männlichen Sprecher; eine Sprecherin verwendet „ráda“.

\smallskip
\textbf{Transkriptionskonvention.} Die IPA-Spalte ist eine breite, manuell erstellte Aussprachehilfe
für Standardtschechisch, keine enge Transkription einer Aufnahme. Die Lautzeichen richten sich nach
den IPA-Übersichten des Phonetischen Instituts der Karls-Universität. Kurze i/y-Vokale werden als \ipa{ɪ},
lange í/ý als \ipa{iː} geschrieben. Die Länge ist von der Betonung unabhängig; auch unbetonte Vokale
können lang sein. Betonung liegt gewöhnlich am Anfang einer Betonungsgruppe, bei einsilbigen Präpositionen
häufig auf der Präposition. Die Tabelle zeigt eine mögliche Lesart; Satzmelodie und Nebenbetonungen
werden nicht vollständig angegeben. Der Punkt in \ipa{ˈpo.uʒiːvaːm} trennt die Vokalsilben.

\smallskip
\textbf{Besondere Konsonanten.} \ipa{r̝} ist das stimmhafte ř, \ipa{r̝̊} seine stimmlose Variante,
beispielsweise nach p oder t. \ipa{ɦ} steht für das stimmhafte h, \ipa{x} für ch.
\ipa{c, ɟ, ɲ} bezeichnen die palatalen Laute ť, ď, ň, keine deutschen Buchstabennamen.
\ipa{r̩, l̩, m̩} sind silbische Konsonanten, etwa in „prší“, „mlha“ und „sedm“.
Affrikaten sind platzsparend als \ipa{ts, tʃ} ohne Bindebögen geschrieben.
Auslautverhärtung und typische Angleichungen innerhalb von Wörtern sind berücksichtigt;
Angleichungen über Wortgrenzen und Glottisschläge werden nicht durchgehend dargestellt.
„Jsem“ und „jsou“ sind in der häufigen Aussprache ohne anlautendes j wiedergegeben.

\smallskip
\textbf{Namen, Zahlen und Kürzel.} Peter und Sittingbourne stehen mit englischer Aussprache,
„Doverem“ mit einer tschechisch angepassten Form. Icom verwendet eine englische Aussprachevariante,
Loop eine englische Form; Hellschreiber und Liechtenstein sind deutsch transkribiert.
Die übrigen Namen werden tschechisch beziehungsweise tschechisch angepasst gelesen.
Rufzeichen und die meisten Kürzel stehen mit tschechischen Buchstabennamen, nicht im internationalen
Buchstabieralphabet: Q „kvé“, W „dvojité vé“, J „jé“. FISTS wird als Wort gelesen;
das e in eQSL und EHam als englisches \ipa{iː}. Zahlen werden ausgeschrieben:
1960 „tisíc devět set šedesát“, 706 „sedm set šest“, 599 „pět devět devět“ und 1234
„tisíc dvě stě třicet čtyři“. 3.558 MHz wird als Dezimalfrequenz „tři celá pět pět osm megahertzů“ gelesen.
5 Watech wird nach Korrektur als „pěti wattech“, 80m als „osmdesát metrů“ und 20 cm als
„dvacet centimetrů“ gesprochen. Platzhalter \texttt{<\#\#\#\#\#>} erscheinen als Auslassung
und werden nicht mitgesprochen.

\smallskip
\textbf{Fachquellen.} Fonetický ústav, Univerzita Karlova: „Transkripce češtiny pomocí IPA“ und
„Sound Patterns of Czech“:
\url{https://fonetika.ff.cuni.cz/o-fonetice/foneticka-transkripce/transkripce-cj-ipa/};
\url{https://fonetika.ff.cuni.cz/en/czech-phonetics/}.
MZV ČR: „Česká arktická věda spojuje síly a otevírá příležitosti pro vědecko-ekonomickou diplomacii“
(2026), Bezeichnung „Špicberky (Svalbard)“.

\smallskip
\textbf{LaTeX.} UTF-8-Datei; zweimal mit XeLaTeX kompilieren. Benötigt werden DejaVu Sans und die
üblichen LaTeX-Pakete der Präambel. IPA ist direkt als Unicode eingegeben; keine externen Dateien erforderlich.
\end{document}
